\documentclass[11pt,a4paper]{article}
\usepackage[UTF8]{ctex}
\usepackage{geometry}
\geometry{left=2.5cm,right=2.5cm,top=2.5cm,bottom=2.5cm}
\usepackage{amsmath,amssymb,amsthm}
\usepackage{hyperref}
\usepackage{booktabs}
\usepackage{enumitem}
\usepackage{xltabular}
\hypersetup{colorlinks=true,linkcolor=blue,urlcolor=blue}

\newtheorem{definition}{定义}[section]
\newtheorem{proposition}{命题}[section]
\newtheorem{remark}{注记}[section]

\title{安全参数自适应调节的形式化与归档\\
——分析能力四层最优分配}
\author{李龙胜}
\date{\today}

\begin{document}

\maketitle

\begin{center}
\textit{
安全参数的旋钮不应由人手拧动。\\
当攻击方探测强度上升时，系统自动缩短扰动周期、压低可观测差异；\\
当探测强度回落时，系统自动释放安全资源。\\
这是分析能力二阶最优分配的第四个维度。\\
贪心择优统一裁定四层之间的最优比例。
}
\end{center}

\vspace{5mm}

\begin{abstract}
在密码学融合的安全运营蓝图中，安全参数（扰动周期 \(T_{\text{rot}}\)、
可观测差异 \(\Delta\)、伪造强度 \(q_{01}, q_{10}\)）的调节
此前依赖人工触发（攻防演练或事件复盘）。
本文将安全参数调节纳入分析能力的统一分配框架，
作为继告警分析、参数重标定、对手行为分析之后的第四层。
系统通过对手行为分析感知攻击方探测强度的变化，
自动计算安全参数调节的边际收益，
由贪心择优在四层之间动态分配有限的分析能力。
\end{abstract}

\tableofcontents

\section{从三层到四层}

在先前的分析能力二阶最优分配框架中，
防御方总分析能力 \(C\) 被分配至三个层面：
告警分析（\(r_{\text{alert}}\)）、
参数重标定（\(r_{\text{param}}\)）、
对手行为分析（\(r_{\text{opponent}}\)）。
三者满足 \(r_{\text{alert}} + r_{\text{param}} + r_{\text{opponent}} \le 1\)。

当密码学组件（同态加密、零知识证明、安全多方计算）嵌入系统后，
安全参数的自适应调节成为一个独立的能力消耗源：
缩短扰动周期需要更频繁的台账解密与重加密，
降低可观测差异需要更高精度的参数控制，
增强伪造响应需要消耗额外的分析能力来生成伪造操作。
这些都不是无代价的。
因此，安全参数调节应被视为第四个分析能力消费者。

\section{四层分配的贪心择优}

防御方在每个周期比较四个层面的边际收益和边际成本，
由贪心择优统一裁定：
\[
\frac{\partial U_{\text{alert}}}{\partial r_{\text{alert}}} =
\frac{\partial U_{\text{param}}}{\partial r_{\text{param}}} =
\frac{\partial U_{\text{opponent}}}{\partial r_{\text{opponent}}} =
\frac{\partial U_{\text{sec}}}{\partial r_{\text{sec}}}
\]
其中 \(U_{\text{sec}}\) 是安全参数调节的净收益——
压缩攻击方探测效率所带来的漏报损失减少，
减去调节操作本身消耗的分析能力。

\section{自适应触发逻辑}

安全参数调节的边际收益并非恒定，
它依赖于攻击方当前的探测强度。
对手行为分析（第3层）持续监测
\(\Delta N\)（新源IP涌现率）、探测密度异常等指标。
当这些指标超过预设阈值时，
系统判断攻击方探测强度上升，
\(r_{\text{sec}}\) 的边际收益自动升高，
贪心择优从其他层面抽调分析能力给安全参数调节：
\begin{itemize}
    \item 缩短扰动周期 \(T_{\text{rot}}\)。
    \item 压低可观测差异 \(\Delta\)。
    \item 增强伪造强度 \(q_{01}, q_{10}\)。
\end{itemize}
当探测强度回落并持续一段时间后，
系统自动降级——恢复较长的 \(T_{\text{rot}}\) 和较宽松的 \(\Delta\)，
释放分析能力给告警分析和参数重标定。

\section{贪心最优性的保证}

贪心最优性定理（生成步系统第一卷定理1）保证：
在成本函数的凸性和进展单调性条件下，
每步贪心择优等价于全局最优。
安全参数调节的成本函数（扰动算力消耗）
对安全级别是严格凸的——
每提升一级安全级别所需的边际算力递增。
因此，四层贪心分配在动态变化中自动趋近全局最优，
无需全局规划。

\section{与已有体系的对接}

本方案不改变四卷体系中的任何已有结论。
它仅将第4卷产品蓝图中“安全参数由人工设定”的
隐含假设显式化，并替换为自适应调节机制。
朋友审查中识别的第17条问题（扰动周期时机）
在自适应框架中得到彻底解决——扰动频率和幅度
不再以事件为锚，而是由系统根据实时探测强度动态决定。

\section{诚实标注}

\begin{center}
\begin{xltabular}{\textwidth}{c|X|c}
\toprule
\textbf{命题} & \textbf{状态} & \textbf{层级} \\
\midrule
四层分析能力分配框架 & 结构完整，与已有三层框架无缝衔接 & II \\
自适应触发逻辑 & 基于已有饱和感知指标的扩展，定性明确 & II \\
贪心择优等价于全局最优 & 凸性条件已验证，定理骨架已证 & II \\
具体阈值标定 & 需在部署环境中建立探测强度基线 & III \\
\bottomrule
\end{xltabular}
\end{center}

\section{结语}

安全参数的旋钮不再由人手拧动。
它被接入了分析能力二阶最优分配的同一套贪心择优回路。
攻击方探测升温时系统自动变硬，
攻击方蛰伏时系统自动节能。
安全运营系统从此具备了自适应安全级别的能力。

\vspace{5mm}
\noindent\textbf{文档信息}：
\begin{itemize}
    \item 作者：李龙胜
    \item 版本：第一版（安全参数自适应调节归档）
    \item 许可：CC BY 4.0
    \item 前置依赖：四卷体系、外部审查修正与补充、密码学融合蓝图
\end{itemize}

\end{document}